\subsection{Nucleo ricorsivo della rete}

La sequenza di token prodotta dall'encoder viene elaborata dal nucleo ricorsivo della TRM, che costituisce la componente centrale del modello. In questa fase è stato mantenuto il principio generale di funzionamento dell'architettura originale, intervenendo principalmente sull'integrazione delle rappresentazioni prodotte dall'encoder e sulla configurazione del processo ricorsivo.

\subsubsection{Rappresentazioni latenti}

Il nucleo mantiene due stati interni, indicati con $z_H$ e $z_L$, entrambi organizzati come sequenze di token. Considerando una dimensione nascosta pari a 256, ciascuno presenta la forma:

\[
    [B, N^2, 256]
\]

dove $B$ indica la dimensione del batch e $N$ quella della scacchiera.

I due stati vengono aggiornati iterativamente attraverso un blocco di ragionamento condiviso. Durante ogni ciclo, $z_L$ viene aggiornato utilizzando le informazioni contenute in $z_H$ e la rappresentazione prodotta dall'encoder; successivamente, il nuovo $z_L$ viene utilizzato per aggiornare $z_H$.

Una scelta specifica dell'adattamento al Go riguarda l'integrazione delle informazioni dell'encoder durante l'intero processo ricorsivo. La rappresentazione iniziale della scacchiera viene infatti reinserita a ogni aggiornamento di $z_L$, anziché essere utilizzata esclusivamente per inizializzare il modello:

\[
    z_L \leftarrow f(z_L, z_H + x)
\]

dove $x$ rappresenta la sequenza di embedding prodotta dal \texttt{GoEncoder} e $f$ indica il blocco di ragionamento.

Questa scelta permette di mantenere disponibili le informazioni relative allo stato originale della partita durante le iterazioni. Al termine del processo, lo stato $z_H$ viene utilizzato come rappresentazione finale delle intersezioni e passato alla successiva componente del modello.

\subsubsection{Blocco di ragionamento}

L'aggiornamento degli stati latenti è affidato a un blocco di ragionamento basato su self-attention e su una componente feed-forward. Il blocco condivide i propri parametri tra le diverse applicazioni, consentendo di riutilizzare la stessa trasformazione durante il processo ricorsivo.

La configurazione adottata prevede una dimensione nascosta di 256, quattro teste di attenzione da 64 dimensioni ciascuna e una componente feed-forward basata su SwiGLU, con fattore di espansione pari a 2. La normalizzazione viene effettuata mediante RMSNorm, mentre l'attenzione causale è disabilitata.

Quest'ultima scelta deriva dalla natura della sequenza elaborata: i token rappresentano le intersezioni di uno stato completo della scacchiera e non una sequenza temporale ordinata.
Ogni posizione può quindi contribuire all'elaborazione delle altre, senza imporre vincoli causali.

Le connessioni residuali e la normalizzazione contribuiscono a mantenere stabili le rappresentazioni durante i successivi aggiornamenti ricorsivi.

\subsubsection{Parametrizzazione e processo ricorsivo}

La configurazione adottata per il nucleo ricorsivo prevede i seguenti parametri:

\begin{table}[H]
    \centering
    \begin{tabular}{l|c}
        \textbf{Parametro} & \textbf{Valore} \\
        \hline
        \texttt{H\_cycles} & 2               \\
        \texttt{L\_cycles} & 2               \\
        \texttt{L\_layers} & 1               \\
        \texttt{H\_layers} & Non utilizzato
    \end{tabular}
    \caption{Configurazione del processo ricorsivo.}
    \label{tab:recursive_core_config}
\end{table}

Una scelta progettuale rilevante riguarda la condivisione dei parametri tra gli aggiornamenti dei due stati latenti. Viene infatti utilizzato un unico modulo \texttt{L\_level}, anziché istanziare componenti distinte per l'elaborazione di $z_L$ e $z_H$.

Con la configurazione adottata, ciascun ciclo $H$ comprende due aggiornamenti consecutivi di $z_L$, seguiti da un aggiornamento di $z_H$. Il processo viene ripetuto per due cicli, producendo complessivamente quattro aggiornamenti di $z_L$ e due di $z_H$.

Anche la gestione del gradiente segue una scelta specifica: il primo ciclo $H$ viene eseguito senza mantenere il grafo computazionale necessario alla backpropagation, mentre il secondo viene eseguito con il calcolo del gradiente attivo. In questo modo è possibile effettuare più aggiornamenti ricorsivi contenendo il consumo di memoria durante l'addestramento.

A differenza della versione originale della TRM, in questa implementazione non viene utilizzato un meccanismo di arresto adattivo basato su una Q-head. Il numero di iterazioni è determinato direttamente dai parametri di configurazione e rimane fisso durante l'elaborazione.

Al termine del processo, il nucleo restituisce lo stato $z_H$, che costituisce l'ingresso della Policy Head incaricata di produrre i punteggi associati alle possibili mosse.
